\section{Method and Governing Equations}

All terms are referred to the radio connector plane.  Power quantities are in dBm;
gains and losses are in dB.  For frequency in GHz and distance in km, the basic
free-space transmission loss is
\begin{equation}
  L_{fs}=92.45+20\log_{10}(f_{\mathrm{GHz}})
               +20\log_{10}(d_{\mathrm{km}}).
  \label{eq:fspl}
\end{equation}
The equivalent isotropically radiated power and clear-air receive level are
\begin{align}
  \mathrm{EIRP} &= P_t-L_{tf}-L_{tc}+G_t, \\
  P_r &= \mathrm{EIRP}-L_{fs}-A_g-L_p-L_o+G_r-L_{rf}-L_{rc}.
  \label{eq:pr}
\end{align}
The receiver noise floor, planning threshold, and fade margin are
\begin{align}
  N &= -174+10\log_{10}(B_{\mathrm{Hz}})+NF, \\
  P_{th} &= N+(C/N)_{req}+L_i, \\
  M_f &= P_r-P_{th}.
  \label{eq:margin}
\end{align}
The first-Fresnel-zone radius at a point on the path is
\begin{equation}
  r_1=\sqrt{\frac{\lambda d_1d_2}{d_1+d_2}},
\end{equation}
with at least 60\% clearance used here as a preliminary screening criterion.

\begin{summarybox}
\textbf{Reference basis.} Equation~\eqref{eq:fspl} follows ITU--R P.525-5.
Propagation screening and availability assessment should follow ITU--R P.530-19;
specific rain attenuation uses ITU--R P.838-3.  Apply the editions required by the
project and the national spectrum authority.
\end{summarybox}
